% Chapter 17 — Scope. Author: KLEPPMANN. Budget 3 pp (G4: +1 pp for the open-problems index).
% Discharges the Constitution v1.1 "Scope" section.
% Hosts the open-problems index and the constitutional-parking rule (G4).
% Target of Chapter~\ref{ch:requirements} V1 (model retention outside scope) via \label{ch:scope}.
% No thread episodes: a scope chapter cites; it does not re-run threads.
\section{Scope}\label{ch:scope}

This chapter discharges the \emph{Scope} section of the constitution.

Scope fixes the boundary of the specification: what the specification covers, and what it
reconciles against at its perimeter but never performs. The constitution states the
in-scope object in one sentence --- the recording, valuation, and lifecycle of
trading-book positions held at fair value, and the interface that projects committed
activity into settlement. Everything below is that sentence made exact, component by
component, together with the explicit list of external authorities that lie outside it.

\subsection{In scope}

The specification covers the following components, each defined and proved in the chapter
named. Nothing in scope lives outside a chapter; nothing below is left to an implementation
to invent.

\begin{itemize}[leftmargin=*, itemsep=2.5pt]
\item \textbf{The map-then-fold architecture and the projection machinery} that derives
  every view from the log --- Chapter~\ref{ch:picture}.
\item \textbf{The immutable log and the objects} --- units, wallets, balances, the atomic
  move, the transaction, and the signed coordinate vector \emph{(owned, lent, posted)} ---
  Chapter~\ref{ch:objects}.
\item \textbf{The three machines} --- the Event Monitor, the Events Executor, and the
  Transaction Executor --- and the single-writer door --- Chapter~\ref{ch:machines}.
\item \textbf{The smart-contract framework}, its event-triggering discipline, and the
  purity and idempotence requirements --- Chapter~\ref{ch:contracts}.
\item \textbf{The three homes} --- ProductTerms, UnitStatus, PositionState --- and the
  conditions that render illegal states unrepresentable --- Chapter~\ref{ch:homes}.
\item \textbf{Valuation and profit and loss} --- NAV as a projection, path-independent PnL,
  and deposit neutrality --- Chapter~\ref{ch:valuation}.
\item \textbf{Ingestion, corporate actions, and the market data operator} --- the recorded
  observation and the operator intrinsic to each (data-kind, event-kind) pair ---
  Chapter~\ref{ch:marketdata}.
\item \textbf{Collateral, margin, and lending} --- the regime key, the coverage invariant,
  and the determination/payment split --- Chapter~\ref{ch:collateral}.
\item \textbf{Virtual ledgers, strategies, and the total return swap} ---
  Chapter~\ref{ch:virtual}.
\item \textbf{The settlement interface} --- the projection of committed activity into
  settlement --- Chapter~\ref{ch:settlement}.
\item \textbf{Presentation and reporting projections} --- every reporting figure read from
  coordinates, units, and declared terms alone --- Chapter~\ref{ch:reporting}.
\item \textbf{Alignment with the Common Domain Model} --- shown, never adopted as authority
  --- Chapter~\ref{ch:cdm}.
\item \textbf{The invariant catalogue} and its executable properties ---
  Chapter~\ref{ch:invariants}.
\item \textbf{Testability and the executable-check regime} ---
  Chapter~\ref{ch:testability}.
\item \textbf{The minimum requirements} any conformant implementation must meet ---
  Chapter~\ref{ch:requirements}.
\end{itemize}

\subsection{Out of scope}

The following are external authorities. The ledger reconciles against each at its boundary;
it performs none of them. Each is a function of some other system, and the ledger's
closure holds precisely because it does not absorb them.

\begin{itemize}[leftmargin=*, itemsep=2.5pt]
\item \textbf{Price formation} --- price discovery and executable pricing. The ledger
  records the outputs of pricing models and never runs one (Chapter~\ref{ch:valuation};
  Chapter~\ref{ch:requirements}, V1). The discovery of price, and the model that produces
  it, are external.
\item \textbf{Legal agreements} --- the enforceability and validity of contracts. Every
  right and obligation a legal contract creates is represented in the ledger as a unit
  (Chapter~\ref{ch:contracts}), but the legal standing of the instrument is decided
  outside the boundary.
\item \textbf{The settlement infrastructure itself} --- settlement finality, payment-system
  access, and custodian and CSD connectivity. The ledger says \emph{what} settles; the
  infrastructure decides how, and whether (Chapter~\ref{ch:settlement}).
\item \textbf{Regulatory submission} --- submission gateways and report formatting. Every
  report is a projection of the record (Chapter~\ref{ch:reporting}); its transmission and
  its per-regime format are external.
\item \textbf{Reference-data authority} --- the mastering of instrument, counterparty, and
  calendar reference data. The ledger reconciles against it at the boundary; it does not
  master it.
\item \textbf{Model governance and retention} --- the validation, approval, and retention
  of pricing and strategy models, and accounting-policy decisions. Reproducing a valuation
  from recorded inputs and identified prices is unconditional; reproducing a model's number
  requires the model itself, whose retention is a governance matter outside the ledger's
  scope (the anchor of Chapter~\ref{ch:requirements}, V1).
\end{itemize}

Three companion documents are named in the Exclusions Register as out-of-scope
deliverables --- expansions the specification points to but does not contain, and not
functions the ledger performs: the \emph{Worked-Examples Volume} (register lines E71 and
E72), holding the dividend-forecast, special-dividend, spin-off, and elective-merger
worked examples; the \emph{SBL Operations Companion} (register line E73), holding the full
securities-lending operational lifecycle --- locates, buy-ins, SFTR and CSDR discipline,
cascade recalls; and the \emph{Track B graphs-interpreter companion} (register line E74),
holding the generic term-sheet interpreter, the CDM round-trip, and term-sheet rendering.
Each is named so its absence from this specification is deliberate and auditable, never
silent.

\subsection{The open-problems index and the constitutional-parking rule}

Two registers live in this chapter. The first names the design questions the specification
leaves open; the second parks any genuine conflict with the constitution.

\paragraph{The open-problems index.} The constitution requires that implementation
questions be named and left, not silently dropped. Each open design question is named here
and carried in full in the \emph{Open-Problems and Escalation Register} companion. The
top open risk is close-out and netting algebra: the collateral ruling made
per-netting-set claim units first-class (Chapter~\ref{ch:collateral}) precisely so that
close-out is representable, yet the algebra that nets and closes those units out is not yet
designed. It entails, on a counterparty default, valuing each claim at close-out --- each
priced by the taker --- netting the per-netting-set claims and obligations to a single
figure, and decoupling recovery from that figure; it is named here so it cannot be
improvised. Named alongside it: managed-account fee accrual and NAV attribution, which the
constitution's \S6 defers to the detailed specification and which this specification
therefore names but does not work; the correction
algebra beyond the single compensating transaction; inter-ledger federation and
eliminations; the fails-cascade obligation's full valuation and lifecycle --- its mint by
the settlement-fail contract, its valuation as a unit, its resolution or supervised
write-off, and the accounting of its paired delivery-versus-payment cash leg
(Chapter~\ref{ch:collateral}); and the remaining items of the companion register. Naming is the whole
obligation; none of these is a gap in what the specification does claim.

\paragraph{The constitutional-parking rule.} The second register governs conflict with
Constitution v1.1 itself.

\begin{principle}[Constitutional parking]
A genuine conflict between this specification and the constitution --- discovered in
drafting or in review --- is parked in this index with the exact proposed amendment text.
It is never silently amended and never fudged. Constitutional amendments are the owner's
alone, under the constitution's closing \emph{Authority and Amendment} section; the
one-way direction of authority forbids an author from resolving a conflict by editing the
constitution. Parking makes the conflict visible and hands the owner a precise, decidable
choice.
\end{principle}

\emph{Current status: the index of parked amendments is empty.} The drafting of
Chapters~\ref{ch:objective} through~\ref{ch:requirements} surfaced no conflict with the
constitution; every author's completion report filed none, and the drafting of this
chapter surfaced none. The index is empty but live: the mechanism stands, numbered and
in place, so that any conflict a later revision or review discovers has a home and a fixed
procedure rather than an ad hoc one.

\subsection{Authority}

The direction of authority is one-way. This specification is rated against the
constitution, and against nothing else; the constitution is rated against nothing. Every
chapter opened by naming the sections it discharges, and this one discharges the
constitution's Scope section. Where a future improvement requires departing from the
constitution, the constitution is amended first --- through the parking rule above and the
owner's decision --- never the specification bent in silence to fit it, and never the
constitution bent to fit the specification.
